% !TEX root = ../main.tex
%--- Slide 3: il modello di rumore (circa 1 min 15 s) ------------------------------------
\begin{frame}{Il modello di rumore}
  \small
  \begin{itemize}\setlength{\itemsep}{2pt}
    \item \textbf{White noise}: errori di misura indipendenti da un TOA all'altro (fissato)
    \item \textbf{Rumori rossi}: processi stocastici a legge di potenza come il GWB ($A$, $\gamma$):
          potenza alle basse frequenze, cioè variazioni lente, su scale di anni
  \end{itemize}
  \vspace{2mm}
  \centering
  \renewcommand{\arraystretch}{1.1}
  \begin{tabular}{lllc}
    \toprule
    Processo & Origine & Tipo & Ritardo $\propto$ \\
    \midrule
    GWB                       & onde gravitazionali   & acromatico & $\nu^{0}$ \\
    Red noise intrinseco (RN) & rotazione della pulsar & acromatico & $\nu^{0}$ \\
    Variazioni di DM          & elettroni liberi & \textbf{cromatico} & $\nu^{-2}$ \\
    Variazioni di scattering (SV) & mezzo interstellare & \textbf{cromatico} & $\nu^{-4}$ \\
    \bottomrule
  \end{tabular}
  \vspace{2mm}
  \begin{itemize}\setlength{\itemsep}{2pt}
    \item \textbf{Cromatico}: il ritardo cresce a bassa frequenza radio $\nu$:
          per separarli servono osservazioni a più frequenze
    \item GWB e RN hanno la stessa firma acromatica: in una singola pulsar non si distinguono
    \item Modello di riferimento: RN + DM + SV, \textbf{6 parametri liberi}
  \end{itemize}
\end{frame}
